\setcounterpageref{ex}{0}
\TaoPhieuKiemTra{12}

\setcounter{ex}{0}
\setcounter{bt}{0}

\textbf{\large I. PHẦN TRẮC NGHIỆM (3,0 điểm)}

{\itshape Khoanh tròn vào chữ cái đặt trước câu trả lời đúng. Mỗi câu đúng được 0,5 điểm.}

\begin{ex} % Câu 1
Số bé nhất có hai chữ số giống nhau là
\choice
{$10$}
{$22$}
{$00$}
{\True $11$}
\loigiai{
Số bé nhất có hai chữ số giống nhau là $11$.\\
\textbf{Đáp án đúng: D}
}
\end{ex}

\begin{ex} % Câu 2
Số bị trừ trong phép tính $94 - 40 = 54$ là
\choice
{$40$}
{$54$}
{\True $94$}
{$134$}
\loigiai{
Số bị trừ là $94$.\\
\textbf{Đáp án đúng: C}
}
\end{ex}

\begin{ex} % Câu 3
Cái bút chì của em dài khoảng
\choice
{\True $15\text{ cm}$}
{$15\text{ kg}$}
{$15\text{ }l$}
{$15\text{ dm}$}
\loigiai{
Chiều dài bút chì đo bằng xăng-ti-mét ($15\text{ cm}$).\\
\textbf{Đáp án đúng: A}
}
\end{ex}

\begin{ex} % Câu 4
Hình chữ nhật có tất cả bao nhiêu đỉnh?
\choice
{$3$ đỉnh}
{\True $4$ đỉnh}
{$5$ đỉnh}
{$6$ đỉnh}
\loigiai{
Hình chữ nhật có $4$ đỉnh và $4$ cạnh.\\
\textbf{Đáp án đúng: B}
}
\end{ex}

\begin{ex} % Câu 5
Điền dấu thích hợp: $36 + 12 \dots 50$.
\choice
{$>$}
{$=$}
{\True $<$}
{Không so sánh được}
\loigiai{
$36 + 12 = 48$. Vì $48 < 50$ nên điền dấu $<$.\\
\textbf{Đáp án đúng: C}
}
\end{ex}

\begin{ex} % Câu 6
Tính: $11 - 5 + 8 = \dots$
\choice
{\True $14$}
{$13$}
{$15$}
{$12$}
\loigiai{
$11 - 5 = 6$; $6 + 8 = 14$.\\
\textbf{Đáp án đúng: A}
}
\end{ex}

\vspace{4mm}
\textbf{\large II. PHẦN TỰ LUẬN (7,0 điểm)}

{\itshape Thí sinh trình bày chi tiết bài làm vào giấy kiểm tra.}

\begin{bt}[1,0 điểm]
Tính nhẩm:
\begin{multicols}{2}
\begin{enumerate}[a)]
  \item $8 + 6$
  \item $7 + 7$
  \item $14 - 5$
  \item $16 - 6$
\end{enumerate}
\end{multicols}
\loigiai{
\begin{itemize}
  \item $8 + 6 = 14$
  \item $7 + 7 = 14$
  \item $14 - 5 = 9$
  \item $16 - 6 = 10$
\end{itemize}
}
\end{bt}

\begin{bt}[1,5 điểm]
Đặt tính rồi tính:
\begin{multicols}{3}
\begin{enumerate}[a)]
  \item $47 + 43$
  \item $79 - 37$
  \item $65 + 52$
\end{enumerate}
\end{multicols}
\loigiai{
\begin{itemize}
  \item $47 + 43 = 90$
  \item $79 - 37 = 42$
  \item $65 + 52 = 117$
\end{itemize}
}
\end{bt}

\begin{bt}[1,0 điểm]
Tính giá trị biểu thức kèm đơn vị đo:
\begin{enumerate}[a)]
  \item $28\text{ kg} + 7\text{ kg} - 9\text{ kg}$
  \item $37\text{ }l - 13\text{ }l + 6\text{ }l$
\end{enumerate}
\loigiai{
\begin{enumerate}[a)]
  \item $28\text{ kg} + 7\text{ kg} - 9\text{ kg} = 35\text{ kg} - 9\text{ kg} = 26\text{ kg}$.
  \item $37\text{ }l - 13\text{ }l + 6\text{ }l = 24\text{ }l + 6\text{ }l = 30\text{ }l$.
\end{enumerate}
}
\end{bt}

\begin{bt}[1,5 điểm]
Bạn Nam có $27$ con tem. Bạn Huy có ít hơn bạn Nam $4$ con tem. Hỏi bạn Huy có bao nhiêu con tem?
\loigiai{
Số con tem bạn Huy có là:
\[ 27 - 4 = 23 \text{ (con tem)} \]
\textbf{Đáp số:} $23$ con tem.
}
\end{bt}

\begin{bt}[1,0 điểm]
Một thùng chứa đầy $47\text{ }l$ dầu ăn. Buổi sáng cửa hàng đã bán ra $16\text{ }l$. Hỏi trong thùng còn lại bao nhiêu lít dầu ăn?
\loigiai{
Trong thùng còn lại số lít dầu ăn là:
\[ 47 - 16 = 31\text{ }l \]
\textbf{Đáp số:} $31\text{ }l$ dầu ăn.
}
\end{bt}

\begin{bt}[1,0 điểm]
Hình vẽ có đoạn thẳng $AB$. Nếu lấy thêm hai điểm $C$ và $D$ nằm giữa $A$ và $B$ thì có tất cả bao nhiêu đoạn thẳng?
\loigiai{
Các điểm thẳng hàng theo thứ tự là $A, C, D, B$.\\
Các đoạn thẳng tạo thành gồm:
- Đoạn đơn: $AC, CD, DB$ (3 đoạn).
- Đoạn đôi: $AD, CB$ (2 đoạn).
- Đoạn ba: $AB$ (1 đoạn).\\
Tổng cộng: $3 + 2 + 1 = 6$ đoạn thẳng.\\
\textbf{Đáp số:} $6$ đoạn thẳng.
}
\end{bt}
